% SPDX-License-Identifier: GPL-3.0-only
% Licensed under GNU GPL version 3; https://www.gnu.org/licenses/gpl-3.0.html . See LICENSE at repository root.
\section{The supported sequent rules}\label{sec:rules}

\subsection{How to read the catalog}
Every displayed inference goes from the premises above the line to the
conclusion below it. Every source and premise must be well formed. A display
with several premises requires all of them; it does not assert that they
can all be recovered from a valid conclusion. This distinction matters for
the implemented left box-choice and left box-star rules.

To specify ordered context replacement compactly, fix context splits and put
\begin{align*}
 R(F)&=(\seq{\Gamma}{\Delta_0,F,\Delta_1}),\\
 L(F)&=(\seq{\Gamma_0,F,\Gamma_1}{\Delta}).
\end{align*}
Here $F$ is a labelled occurrence, including a tag when present.
$R(F_1,\ldots,F_n)$ or $L(F_1,\ldots,F_n)$ replaces the principal by that
ordered list. $R()$ and $L()$ remove it. In the tables below,
``$[F_1,F_2];[F_3]$'' means two premises, with replacement lists
$[F_1,F_2]$ and $[F_3]$ in that order. It never means one premise with three
occurrences. Tags are suppressed in the logical catalog; their routing is
specified in Section~\ref{sec:progress}.
\sources{\rocq{Rules.v}{right_targets}, \rocq{Rules.v}{left_targets}}

\subsection{Axioms and manual first-order leaves}
The structural axiom shapes are
\[
 \Rule{Id}{}{\seq{\Gamma,\ann{\ell}{A}}{\ann{\ell}{A},\Delta}}
 \qquad
 \Rule{True.R}{}{\seq{\Gamma}{\ann{\ell}{\top},\Delta}}
 \qquad
 \Rule{False.L}{}{\seq{\Gamma,\ann{\ell}{\bot}}{\Delta}}.
\]
Identity requires equal labels and formulas; it ignores differences in the
tags of the matched occurrences. The matches can occur anywhere in the
ordered contexts. In addition, a singleton-state antecedent
$\ann{U}{\Suf{A}{B}}$ closes a judgment because a singleton has no proper
nonempty suffix. All these axioms still require $\WF(J)$.
\sources{\rocq{Rules.v}{identity_axiom}, \rocq{Rules.v}{truth_right_axiom},
  \rocq{Rules.v}{falsity_left_axiom}, \rocq{Rules.v}{suffix_singleton_axiom}}

The terminal interface \code{close_atoms} accepts only predicate or equality
occurrences at singleton state-expression labels. The broader
\code{close_pure} accepts any pure formula at those labels. Both require a
caller-supplied Rocq proof of
\[
 \forall w.\ (\forall o\in\Gamma.\ w\sat o)
                 \to(\exists o\in\Delta.\ w\sat o).
\]
The shape check is not a decision procedure for the displayed proposition.
An empty terminal shape, for instance, cannot close the false empty sequent
without an impossible proof.
\sources{\rocq{Rules.v}{atom_terminalb}, \rocq{Rules.v}{pure_terminalb},
  \rocq{Rules.v}{close_atoms}, \rocq{Rules.v}{close_pure}}

\subsection{Boolean rules}
For $F=\ann{\ell}{A\land B}$ or $\ann{\ell}{A\lor B}$, the ordered
replacement premises are
\begin{center}
\begin{tabular}{@{}lll@{}}
\toprule
Principal formula & Right replacement lists & Left replacement lists\\
\midrule
$A\land B$ & $[\ann{\ell}{A}];[\ann{\ell}{B}]$
             & $[\ann{\ell}{A},\ann{\ell}{B}]$\\
$A\lor B$ & $[\ann{\ell}{A},\ann{\ell}{B}]$
             & $[\ann{\ell}{A}];[\ann{\ell}{B}]$\\
\bottomrule
\end{tabular}
\end{center}
\sources{\rocq{Rules.v}{conjunction_right}, \rocq{Rules.v}{conjunction_left},
  \rocq{Rules.v}{disjunction_right}, \rocq{Rules.v}{disjunction_left}}

Negation and implication move some operands across the sequent:
\begin{gather*}
 \Rule{Not.R}{\seq{\Gamma,\ann{\ell}{A}}{\Delta_0,\Delta_1}}
   {\seq{\Gamma}{\Delta_0,\ann{\ell}{\neg A},\Delta_1}},\\[0.8em]
 \Rule{Not.L}{\seq{\Gamma_0,\Gamma_1}{\ann{\ell}{A},\Delta}}
   {\seq{\Gamma_0,\ann{\ell}{\neg A},\Gamma_1}{\Delta}},\\[0.8em]
 \Rule{Impl.R}{\seq{\Gamma,\ann{\ell}{A}}{\Delta_0,\ann{\ell}{B},\Delta_1}}
   {\seq{\Gamma}{\Delta_0,\ann{\ell}{A\to B},\Delta_1}},\\[0.8em]
 \Rule{Impl.L}{\seq{\Gamma_0,\Gamma_1}{\ann{\ell}{A},\Delta}
       \qquad\seq{\Gamma_0,\ann{\ell}{B},\Gamma_1}{\Delta}}
   {\seq{\Gamma_0,\ann{\ell}{A\to B},\Gamma_1}{\Delta}}.
\end{gather*}
\sources{\rocq{Rules.v}{negation_right}, \rocq{Rules.v}{negation_left},
  \rocq{Rules.v}{implication_right}, \rocq{Rules.v}{implication_left}}

Disjunction and implication have named derived rules even though their roots
are primitive negations. These are separate proved local instances.

\subsection{Quantifier rules}
All bodies $P$ below must be pure. Instantiation uses capture-avoiding
substitution and a well-formed term $t$ of the binder's sort:
\[
 \Rule{All.L}{L(\ann{\ell}{P[t/x]})}{L(\ann{\ell}{\forall x.P})}
 \qquad
 \Rule{Ex.R}{R(\ann{\ell}{P[t/x]})}{R(\ann{\ell}{\exists x.P})}.
\]
The sort check remains mandatory for a vacuous binder. Eigenvariable rules are
\[
 \Rule{All.R}{R(\ann{\ell}{P[y/x]})}{R(\ann{\ell}{\forall x.P})}
 \qquad
 \Rule{Ex.L}{L(\ann{\ell}{P[y/x]})}{L(\ann{\ell}{\exists x.P})},
\]
where $y$ has the sort of $x$ and is absent from all surrounding contexts,
the principal label, and the body, according to the conservative name
collection in Section~\ref{sec:sequents}.
\sources{\rocq{Rules.v}{universal_left}, \rocq{Rules.v}{existential_right},
  \rocq{Rules.v}{universal_right}, \rocq{Rules.v}{existential_left}}

Freshness must preserve not only context truth but its semantic ranks.
The eigenvariable counterstep uses an involution exchanging two envelope
values at $y$ in every state. This reversible transformation transports
counterexamples without changing their component lengths.
\sources{\rocq{Semantics.v}{swap_values_involutive},
  \rocq{Semantics.v}{eigen_instance}, \rocq{Progress.v}{minimum_switch}}

\subsection{Assignments and tests}
Assignment rules expose a final singleton state expression. Put
$U'=\upd{x:=t}{U}$ and $\ell=\sigma\jlabel U$. Then
\begin{gather*}
 \Rule{B.Assign.R}{R(\ann{(\sigma\jlabel U)\jlabel U'}{A})}
   {R(\ann{\sigma\jlabel U}{[x:=t]A})},\qquad
 \Rule{B.Assign.L}{L(\ann{(\sigma\jlabel U)\jlabel U'}{A})}
   {L(\ann{\sigma\jlabel U}{[x:=t]A})}.
\end{gather*}
The diamond pair has the same premises, replacing the conclusion's box by
$\langle x:=t\rangle$. The body is not syntactically substituted.
The label retains both the old state and the new state.
\sources{\rocq{Rules.v}{box_assignment_right}, \rocq{Rules.v}{box_assignment_left},
  \rocq{Rules.v}{diamond_assignment_right}, \rocq{Rules.v}{diamond_assignment_left}}

For tests, abbreviate $G=\ann{U\jlabel U}{B}$,
$H=\ann{(\sigma\jlabel U)\jlabel U}{A}$, and
$F=\ann{\sigma\jlabel U}{[B?]A}$. The exact test schemes are
\begin{gather*}
 \Rule{B.Test.R}{\seq{\Gamma,G}{\Delta_0,H,\Delta_1}}
   {\seq{\Gamma}{\Delta_0,F,\Delta_1}},\\[0.8em]
 \Rule{B.Test.L}{\seq{\Gamma_0,H,\Gamma_1}{\Delta}}
   {\seq{\Gamma_0,F,G,\Gamma_1}{\Delta}}.
\end{gather*}
Box-test-left consumes a following guard occurrence in the antecedent; it
is not a general unary replacement of $F$ alone. For
$D=\ann{\sigma\jlabel U}{\langle B?\rangle A}$, the diamond schemes are
\begin{gather*}
 \Rule{D.Test.R}{\seq{\Gamma}{\Delta_0,H,\Delta_1}}
   {\seq{\Gamma,G}{\Delta_0,D,\Delta_1}},\\[0.8em]
 \Rule{D.Test.L}{\seq{\Gamma_0,H,G,\Gamma_1}{\Delta}}
   {\seq{\Gamma_0,D,\Gamma_1}{\Delta}}.
\end{gather*}
In the right diamond rule, $G$ is the final antecedent occurrence.
The test guard is observed on $U\jlabel U$, never on $U$ by default.
\sources{\rocq{Rules.v}{box_test_right}, \rocq{Rules.v}{box_test_left},
  \rocq{Rules.v}{diamond_test_right}, \rocq{Rules.v}{diamond_test_left}}

\subsection{Regular program structure}
The following table specifies ordered replacements at the same label $\ell$.
An entry $A$ stands for the labelled occurrence $\ann{\ell}{A}$.
Each row has its named right and left version.
\begin{center}
\begin{tabular}{@{}lll@{}}
\toprule
Principal & Right replacement lists & Left replacement lists\\
\midrule
$[\alpha;\beta]A$ & $[[\alpha][\beta]A]$ & $[[\alpha][\beta]A]$\\
$[\alpha\cup\beta]A$ & $[[\alpha]A];[[\beta]A]$ & $[[\alpha]A];[[\beta]A]$\\
$[\alpha^*]A$ & $[[\top?]A];[[\alpha][\alpha^*]A]$
  & $[[\top?]A];[[\alpha][\alpha^*]A]$\\
$\langle\alpha;\beta\rangle A$ & $[\langle\alpha\rangle\langle\beta\rangle A]$
  & $[\langle\alpha\rangle\langle\beta\rangle A]$\\
$\langle\alpha\cup\beta\rangle A$ & $[\langle\alpha\rangle A,\langle\beta\rangle A]$
  & $[\langle\alpha\rangle A];[\langle\beta\rangle A]$\\
$\langle\alpha^*\rangle A$ & $[\langle\top?\rangle A,\langle\alpha\rangle\langle\alpha^*\rangle A]$
  & $[\langle\top?\rangle A];[\langle\alpha\rangle\langle\alpha^*\rangle A]$\\
\bottomrule
\end{tabular}
\end{center}
\sources{\rocq{Rules.v}{box_sequence_right}, \rocq{Rules.v}{box_sequence_left},
  \rocq{Rules.v}{box_choice_right}, \rocq{Rules.v}{box_choice_left},
  \rocq{Rules.v}{box_star_right}, \rocq{Rules.v}{box_star_left},
  \rocq{Rules.v}{diamond_sequence_right}, \rocq{Rules.v}{diamond_sequence_left},
  \rocq{Rules.v}{diamond_choice_right}, \rocq{Rules.v}{diamond_choice_left},
  \rocq{Rules.v}{diamond_star_right}, \rocq{Rules.v}{diamond_star_left}}

Both left box-choice and left box-star have two premises. They are sound
inferences in the displayed direction, without being the usual invertible
one-premise antecedent expansion of a conjunction. Star unfolding by itself
advertises no strict progress. Its base branch contains the true test and
therefore a genuine stuttering transition when that test is subsequently
expanded.

\subsection{First-state and suffix rules}
For either side $C\in\{L,R\}$,
\[
 \Rule{First}{C(\ann{U}{A})}{C(\ann{U\jlabel\sigma}{\Fst{A}})}
 \qquad
 \Rule{At}{C(\ann{U}{P})}{C(\ann{U\jlabel\sigma}{P})}.
\]
The body $A$ of \textsc{First} is unrestricted. The \textsc{At} rule requires
$P$ to be a predicate application or equality. The pure-head semantic theorem
is broader, but this named local rule retains the atom boundary.
\sources{\rocq{Rules.v}{first_right}, \rocq{Rules.v}{first_left},
  \rocq{Rules.v}{atom_head_right}, \rocq{Rules.v}{atom_head_left}}

For suffix rules the principal label is $U\jlabel\sigma$, and $\sigma$ must
be nonempty. All replacement formulas below have label $\sigma$:
\begin{center}
\begin{tabular}{@{}lll@{}}
\toprule
Principal & Right replacement lists & Left replacement lists\\
\midrule
$\Suf{A}{B}$ & $[B,A];[B,\Suf{A}{B}]$ & $[B];[A,\Suf{A}{B}]$\\
$\DSuf{A}{B}$ & $[B];[A,\DSuf{A}{B}]$ & $[B,A];[B,\DSuf{A}{B}]$\\
\bottomrule
\end{tabular}
\end{center}
\sources{\rocq{Rules.v}{suffix_right}, \rocq{Rules.v}{suffix_left},
  \rocq{Rules.v}{dual_suffix_right}, \rocq{Rules.v}{dual_suffix_left}}

The order follows the one-step laws~\eqref{eq:suffix-step}
and~\eqref{eq:dual-step}. For suffix-right, the first premise contains no
recursive suffix occurrence. Only the second premise has a principal
continuation that can carry suffix progress.

\subsection{Structural rules and conversions}
Weakening has one context-only premise:
\[
 \Rule{Wk.R}{R()}{R(F)}\qquad\Rule{Wk.L}{L()}{L(F)}.
\]
It is not a zero-premise axiom. Exchange permits explicit permutations;
the positional request API exposes adjacent exchanges. Cut has the exact
context partitions
\[
 \Rule{Cut}{\seq{\Gamma}{F,\Delta}\qquad\seq{\Sigma,F}{\Lambda}}
   {\seq{\Gamma,\Sigma}{\Delta,\Lambda}}.
\]
The displayed cut placement is part of its ordered premise list. Any tag
on $F$ must be legal in both premises; the antecedent requirement therefore
excludes a tagged cut occurrence.
\sources{\rocq{Rules.v}{weakening_right}, \rocq{Rules.v}{weakening_left},
  \rocq{Rules.v}{exchange_rule}, \rocq{Rules.v}{cut_rule}}

World and trace substitutions have one original premise and its
label-substituted conclusion:
\[
 \Rule{Sub.World}{J}{J[U/i]}
 \qquad
 \Rule{Sub.Trace}{J}{J[\sigma/k]}.
\]
The replacements must have valid syntax, and $\sigma$ must be nonempty.
Rank transport uses the corresponding changed semantic world; it is not an
assumption that the original and substituted labels denote the same trace
in one fixed world.
\sources{\rocq{Rules.v}{world_substitution_rule},
  \rocq{Rules.v}{trace_substitution_rule}, \rocq{Progress.v}{rank_label_mapping}}

Label rotation and unit rules have premise $J$ and conclusion obtained by
applying their chosen path operation to every label in $J$. Update rewriting
instead changes exactly one occurrence on one side, with the same formula
and tag, from $\ann{\ell}{A}$ in the conclusion to $\ann{\ell'}{A}$ in the
premise, where the selected update rewrite returns $\ell'$. These conversions
preserve the relevant ranks and advertise no progress.
\sources{\rocq{Rules.v}{label_rotation_rule}, \rocq{Rules.v}{label_unit_rule},
  \rocq{Rules.v}{rewrite_update_right}, \rocq{Rules.v}{rewrite_update_left}}

An exact backlink instance has premise and conclusion $J$. A tag-renaming
backlink has conclusion $J$ and premise $\mathsf{renameTags}(f,J)$, requiring
well-formedness on both sides and hence excluding tag collisions. Both have
empty progress sets. Their use as graph backlinks additionally requires
the graph checks in Section~\ref{sec:certification}.
\sources{\rocq{Rules.v}{exact_backlink}, \rocq{Rules.v}{tag_renaming_backlink}}
